\section*{5 Lower bounds (author source pp.22:10--22:12)}
Lower bounds for depth-$d$ Frege systems for $d>1$, typically follow the following strategy:
\begin{enumerate}\item We first show that the formulas we are trying to refute are robust; namely, after applying a restriction selected at random to them, then with high probability they cannot be refuted with proofs whose lines are, in a certain sense, simple.\item Then we show, through the use of a switching lemma, that applying such a restriction to a short proof will result with high probability in a proof with simple lines.\end{enumerate}
Here we start with $\RRef_{d,N}(F,s)$, which after applying the restrictions will collapse to $\Ref(F,s')$, where $s'$ is polynomially related to $s$. For the part of the overall strategy showing that there cannot be refutations with simple lines, we take, as in[2], simple to mean of small index-width. We say that a variable of the form $D[u,i,b],V[u,i],I[u,j],L[u,v]$ or $R[u,v]$ mentions the index $u$. The index-width of a clause in the variables of $\Ref(F,s)$ is defined as the number of indices mentioned by its variables, and the index-width of a resolution refutation of $\Ref(F,s)$ is the maximum index-width over its clauses. We have:
\begin{originalresult}{Theorem 5.1 [2]}
For all integers $n,s>0$ with $s\le2^n$, and every unsatisfiable CNF $F$ with $n$ variables, every resolution refutation of $\Ref(F,s)$ has index-width at least $s/6n$.
\end{originalresult}
\subsection*{5.1 The robustness of $\RRef_{d,N}$}
We create a distribution on restrictions to the variables of $\RRef_{d,N}(F,s)$ as follows. Suppose $d$ is odd (if $d$ were even, we would exchange the roles of0 and1 in the following construction). For each $S_{d,N}^{\wedge}(x)$ formula in $\RRef_{d,N}(F,s)$, look at its bottom-most $N^{d-1}$ $\bigwedge$ connectives. For each such connective, we decide to ``preserve'' it with probability $1/\sqrt N$, and not to preserve it with probability $1-1/\sqrt N$. For each of the preserved connectives, we leave its first variable unset and set the rest to1. For each variable in the unpreserved connectives, we set it to0 or1 with probability $1/2$ for each choice. The variables of $S_{d,N}^{\vee}(x)$ are set in the same way, except that the set variables of the preserved $\bigvee$ connectives are set to0 instead of1.
Under such restrictions, Sipser functions do not simplify much. For formulas $F$ and $G$, in which each variable appears only once, we say that $F$ contains $G$ if we can get $G$ from $F$ by deleting some of its literals and/or renaming some of its variables.
\begin{originalresult}{Lemma 5.2}
For any $d\ge2$, the probability that $S_{d,N}^{\nu}(x)|_\rho$, where $\nu\in\{\wedge,\vee\}$, does not contain $S_{d-1,N}^{\nu}(x)$ is at most $2^{-\Omega(\sqrt N)}$.
\end{originalresult}
\begin{proof}
We show the lemma for $S_{d,N}^{\wedge}(x)$ and $d$ odd. If $S_{d,N}^{\wedge}(x)|_\rho$ does not contain $S_{d-1,N}^{\wedge}$, then either one of its bottom-most $\bigwedge$ connectives takes the value1, or in one of its depth-2 subformulas, less than $\sqrt N/2$ $\bigwedge$ connectives are preserved. The probability that a bottom-most $\bigwedge$ connective takes the value1 is at most $2^{-\sqrt N/2}$ and the probability that this happens for at least one of the $N^{d-1}$ bottom-most $\bigwedge$ connectives is at most
\[N^{d-1}2^{-\sqrt N/2}\le2^{-\Omega(\sqrt N)}.\]
Now fix a depth-2 subformula $A$ of $S_{d,N}^{\wedge}(x)$. The expected number of preserved $\bigwedge$ connectives in $A|_\rho$ is $N/\sqrt N=\sqrt N$, and by the Chernoff bound, the probability that there are less than $\sqrt N/2$ preserved $\bigwedge$ connectives is at most $2^{-\Omega(\sqrt N)}$. The probability that at least one of the $N^{d-2}$ depth-2 subformulas of $S_{d,N}^{\wedge}(x)$ has less than $\sqrt N/2$ preserved connectives is thus at most
\[N^{d-2}2^{-\Omega(\sqrt N)}\le2^{-\Omega(\sqrt N)}.\]
We conclude that the probability that $S_{d,N}^{\wedge}|_\rho$ does not contain $S_{d-1,N}^{\wedge}$ is at most $2^{-\Omega(\sqrt N)}$.
\end{proof}
\subsection*{5.2 The Furst-Saxe-Sipser switching lemma}
Switching lemmas provide conditions under which a $k$-DNF formula ``switches'' to a $\ell$-CNF formula after applying a restriction created at random. We will use the switching lemma of[8] and a variation tailored for $R^k\Ref(F,s)$ due to[10].
Let $G$ be a $k$-DNF formula over the set of variables $X$. Let $X_1,\ldots,X_r$ be a partition of $X$ into $r$ blocks, and let $\nu\in\bits$. Consider the following distribution over restrictions on $X$: For each block $X_i$, we decide to ``preserve'' $X_i$ with probability $p$, and not to preserve it with probability $1-p$. For each preserved block, we leave one of its variables, say the first in the block, unset, and set all others to $\nu$. For each unpreserved block, we set each of its variables to0 or1 with probability $1/2$ for each value.
We can extract the following lemma from[8,19]. The lemma is implicit in[8,19] with parameters obscured under a big $O$ notation. We present it here in a more general, improved form, with explicit parameters, using decision trees along the lines of[18]. In what follows,$\ln$ denotes the natural logarithm; we preserve the notation $\log$ for the base2 logarithm.
\begin{originalresult}{Lemma 5.3 (see[8,19])}
If $phk2^k\ln N=o(N^{-\varepsilon})$ for some $\varepsilon\in(0,1)$, then
\[\mathbb P[\height(G|_\rho)>kh]\le o(N^{-\varepsilon h})\frac{2^{kh}-1}{2^h-1}.\]
\end{originalresult}
\begin{proof}
The proof is by induction on $k$. If $k=0$, $G$ is a constant and can be represented by a decision tree of height0. Suppose $k>0$. We distinguish between two cases,$G$ being wide and $G$ being narrow. We call $G$ wide if there are at least $h2^k\ln N$ terms in it such that no two of them contain variables from the same block. $G$ is narrow if and only if it is not wide. If $G$ is wide, then
\begin{align*}
\mathbb P[\height(G|_\rho)>kh]&\le\mathbb P[G|_\rho\ne1]\le\left(1-\left(\frac{1-p}{2}\right)^k\right)^{h2^k\ln N}\\
&\le e^{-(1-p)^k h\ln N}=N^{-(1-p)^kh}=o(N^{-\varepsilon h}).
\end{align*}
If $G$ is narrow, then take a maximal set of terms such that no two of them contain variables from the same block, and let $H$ be the set of blocks that contain a variable occurring in some term of this set. $H$ contains at most $hk2^k\ln N$ blocks and every term of $G$ contains some variable (or its negation) from some block in $H$. The probability of the event $A$ that $\rho$ preserves more than $h$ blocks in $H$ is
\[\mathbb P[A]\le\binom{hk2^k\ln N}{h}p^h\le(hk2^k\ln N)^hp^h=o(N^{-\varepsilon h}).\]
Now, let $\pi$ be a restriction that sets the variables of all blocks in $H$, and let $A_\pi$ be the event that $\pi$ is consistent with $\rho$ and $\height((G|_\rho)|_\pi)>(k-1)h$. Notice that $G|_\pi$ is a $(k-1)$-DNF, so by the induction hypothesis,
\[\mathbb P[A_\pi]\le\mathbb P[\height((G|_\pi)|_\rho)>(k-1)h]\le o(N^{-\varepsilon h})\frac{2^{(k-1)h}-1}{2^h-1}.\]
Notice that a restriction $\rho$ that preserves at most $h$ blocks is consistent with at most $2^h$ restrictions $\pi$, so we get
\[\mathbb P\left[A\cup\bigcup_\pi A_\pi\right]\le o(N^{-\varepsilon h})+o(N^{-\varepsilon h})2^h\frac{2^{(k-1)h}-1}{2^h-1}=o(N^{-\varepsilon h})\frac{2^{kh}-1}{2^h-1}.\]
In the event $\left(A\cup\bigcup_\pi A_\pi\right)^c$, i.e. the event that $\rho$ preserves at most $h$ blocks in $H$ and for all restrictions $\pi$ consistent with $\rho$,$\height((G|_\rho)|_\pi)\le(k-1)h$, we can construct a decision tree of height at most $kh$ representing $G|_\rho$ as follows: We query all variables belonging to some block in $H$ left unset by $\rho$ (since $\rho$ preserves at most $h$ blocks in $H$, there are at most $h$ of them), and at each branch $\pi$ of the resulting tree, we append a decision tree of minimum height representing $(G|_\rho)|_\pi$.
\end{proof}
